\documentclass[11pt,a4paper]{article}
\usepackage{fontspec}
\IfFontExistsTF{Montserrat}{\setmainfont{Montserrat}}{}
\usepackage{babel}
\babelprovide[main, import]{dutch}
\usepackage[a4paper,margin=2.4cm]{geometry}
\usepackage{enumitem}
\usepackage{xcolor}
\usepackage{titlesec}
\usepackage{amssymb}
\usepackage{hyperref}

\definecolor{hgteal}{RGB}{22,176,165}
\definecolor{hggrey}{RGB}{120,120,120}

\hypersetup{colorlinks=true, linkcolor=hgteal, urlcolor=hgteal}

\newcommand{\code}[1]{\texttt{\small #1}}

\titleformat{\section}{\Large\bfseries\color{hgteal}}{}{0pt}{}
\titlespacing*{\section}{0pt}{1.4em}{0.5em}

\setlist[itemize]{leftmargin=1.4em, itemsep=0.2em, topsep=0.3em}

\setlength{\parindent}{0pt}
\setlength{\parskip}{0.4em}

\begin{document}

{\Huge\bfseries\color{hgteal} Slide-onderbouwing}\\[0.4em]
{\large Geautomatiseerde security assessment van MCP-servers}\\[0.2em]
{\color{hggrey} Verdediging bachelorproef --- Arthur Van Oudenhove}

\vspace{0.8em}
\hrule
\vspace{0.6em}

Per slide: de tekst op de deck en de onderbouwing uit de bachelorproef
(FINAL-versie: precisie-framing, \'e\'en evaluatiecase). Bronverwijzingen tussen
haakjes zijn de \code{\textbackslash autocite}-keys uit \code{bachproef.bib}.

\section{Slide 2. De poster}
\textbf{Op de deck:} volledige poster, full screen, geen tekst.

\textbf{Achtergrond:} dient als totaalkaart. De navigator licht per slide de
relevante posterzone op.

\section{Slide 3. AI-agents in Fintech, via MCP}
\textbf{Op de deck:}
\begin{itemize}
  \item AI-agents verwerken betalingen en bedrijfskritische taken.
  \item Communicatie loopt via het Model Context Protocol.
  \item Standaardisatie opent aanvalsvlakken: prompt injection, tool poisoning, privilege escalation.
  \item ``56\% van de bedrijven met AI-agents rapporteert MCP-beveiligingsincidenten.''
\end{itemize}

\textbf{Achtergrond (Inleiding, Probleemstelling; Stand van zaken, sectie 1):}
\begin{itemize}
  \item Fintech groeide de afgelopen vijf jaar met 23,58\% per jaar (\code{statista2024fintechgrowth}).
  \item MCP is de de-facto standaard voor agent-naar-systeem-communicatie, vergelijkbaar met de rol van HTTP voor webapplicaties (\code{anthropic2024mcpwhitepaper}).
  \item De drie aanvalsvectoren komen uit (\code{song2025protocolunveilingattackvectors}). Confused-deputy en ``Living off AI'' uit (\code{Cato2025}).
  \item De 56\%-statistiek komt uit een systematische analyse (\code{guo2025systematicanalysismcpsecurity}).
  \item Doelgroep: security teams, DevSecOps-engineers en enterprise architects in gereguleerde sectoren.
\end{itemize}

\section{Slide 4. Onderzoeksvraag}
\textbf{Op de deck:}
\begin{itemize}
  \item Hoofdvraag: hoe kan een framework MCP-servers automatisch testen op kritieke kwetsbaarheden, en hoe verhoudt dekkingsgraad en tijd zich tot een handmatige beoordeling?
  \item Scope: \'e\'en gecontroleerde evaluatiecase (Goat) met een vaste ground truth.
\end{itemize}

\textbf{Achtergrond (Inleiding, Onderzoeksvraag en Doelstelling):}
\begin{itemize}
  \item Deelvragen probleemdomein: welke OWASP-kwetsbaarheden zijn kritiek voor Fintech, welke beperkingen hebben Burp en Semgrep, hoeveel tijd kost een handmatige review en wat zijn de foutbronnen.
  \item Deelvragen oplossingsdomein: recall en precisie, het tijdsverschil, en waar menselijke expertise nodig blijft.
  \item Deliverables: een werkend prototype, een kwantitatieve evaluatie (recall, tijd, precisie) en evidence-based aanbevelingen met expliciete vermelding van de beperkte externe validiteit.
  \item De FINAL-versie spreekt bewust over \'e\'en case. De ``twee cases'' uit de draft is geschrapt.
\end{itemize}

\section{Slide 5. Waarom generieke tools tekortschieten}
\textbf{Op de deck:}
\begin{itemize}
  \item Burp Suite ziet HTTP en JSON-RPC, maar kent de semantiek van MCP-schema's niet.
  \item Semgrep detecteert code-patronen, maar niet de runtime-interactie.
  \item ``Geen contextuele kennis van MCP-specifieke kwetsbaarheden.''
\end{itemize}

\textbf{Achtergrond (Stand van zaken, Status quo en Bestaande tools):}
\begin{itemize}
  \item Een handmatige MCP-review kost gemiddeld 8 tot 16 uur per iteratie (\code{Uproot2025}), wat een bottleneck vormt in agile omgevingen.
  \item Burp onderschept JSON-RPC maar herkent geen kwaadaardige tool-beschrijving. Semgrep mist runtime-analyse.
  \item MCPSafetyScanner (\code{radosevich2025mcpsafetyaudit}) is LLM-agentisch, dus niet-deterministisch, en vereist een externe API (een privacyprobleem in Fintech). Jouw keuze valt op deterministische, vaste detectieregels (zie NF-2 reproduceerbaarheid).
  \item Invariant Labs (\code{Invariant2025}) biedt een losse scanning-tool zonder breder framework.
\end{itemize}

\section{Slide 6. Architectuur}
\textbf{Op de deck:} TikZ-diagram. MCP-server naar Discovery naar SAST en Fuzzing naar Rapportgenerator naar report.json en report.md, binnen een ``Scanner, Docker''-container.

\textbf{Achtergrond (Implementatie; Methodologie, fase 3):}
\begin{itemize}
  \item Drie onafhankelijke modules, sequentieel uitgevoerd. Ze delen \'e\'en \code{Finding}-model (\code{rule\_id}, \code{severity}, \code{tool\_name}, \code{title}, \code{detail}), zodat SAST en fuzzing uniform verwerkt worden.
  \item Discovery roept \code{list\_tools}, \code{list\_prompts} en \code{list\_resources} aan en valideert de respons via Pydantic-modellen.
  \item Gebouwd in Python 3.12 met de offici\"ele \code{mcp}-library. Draait via Docker Compose op een ge\"isoleerd intern netwerk (\code{mcp-lab-net}).
  \item Twee rapporten: report.json (integreerbaar in CI/CD) en report.md (leesbaar, met badges).
\end{itemize}

\section{Slide 7. Statisch en dynamisch}
\textbf{Op de deck:}
\begin{itemize}
  \item SAST (schema): prompt injection in beschrijving, ontbrekende autorisatiecontext, ongebonden parameters.
  \item Fuzzing (runtime): information disclosure, injection-reflectie, autorisatiebypass en tool poisoning.
  \item ``Runtime-gedrag is onzichtbaar in broncode, daarom allebei.''
\end{itemize}

\textbf{Achtergrond (Implementatie, SAST en Fuzzing; Stand van zaken, Geautomatiseerd):}
\begin{itemize}
  \item SAST telt vier regels: SAST-001 prompt injection (regex op \code{ignore rules}, \code{you are now}), SAST-002 ongebonden numeriek, SAST-003 ontbrekende autorisatie, SAST-004 ongebonden string.
  \item Fuzzing telt vier regels: FUZZ-001 information disclosure, FUZZ-002 injection-reflectie, FUZZ-003 autorisatiebypass (plus negatieve bedragen), FUZZ-004 tool poisoning (scant de respons).
  \item Onderbouwing van SAST, DAST en fuzzing: (\code{chess2004static}) en (\code{manes2019art}). De regels staan in losse lijsten, dus een nieuwe regel is een functie toevoegen (NF-1 modulariteit).
\end{itemize}

\section{Slide 8. Demo (screencast)}
\textbf{Op de deck:} still en de aanduiding demo (screencast).

\textbf{Achtergrond (Validatie, Testopstelling):}
\begin{itemize}
  \item Goat bevat vijf opzettelijk ingebedde kwetsbaarheden. V1 prompt injection in de description en V2 geen autorisatie op \code{analyze\_financial\_document}. V3 geen autorisatie en V4 negatieve bedragen op \code{execute\_transfer}. V5 tool poisoning op \code{fetch\_market\_news}.
  \item Volledig gecontaineriseerd op een ge\"isoleerd netwerk, wat een reproduceerbare ground truth geeft.
\end{itemize}

\section{Slide 9. Resultaten op Goat}
\textbf{Op de deck (tabel):} Recall 100\% (5/5) tegenover 80\% (4/5). Precisie minstens 92,3\% (12/13) tegenover 100\% (4/4). Tijd minder dan 30 seconden tegenover 47 minuten. Reproduceerbaarheid gegarandeerd tegenover reviewer-afhankelijk. ``Het tijdsverschil van factor 90 is reviewer-onafhankelijk.''

\textbf{Achtergrond (Validatie, Resultaten en Precisie):}
\begin{itemize}
  \item De scanner produceerde 13 bevindingen: SAST 7 (HIGH 3, MED 1, LOW 3) en fuzzing 6 (HIGH 6). Alle vijf de ground-truth-items werden gedetecteerd, dus recall 100\%.
  \item Precisie is 12/13, ongeveer 92,3\%. De FINAL-versie legt uit dat een klassieke FP-rate (FP gedeeld door FP plus TN) niet zinvol is in een vuln-scan, omdat true negatives niet zinvol te tellen zijn. De debateerbare vinding is FUZZ-002 op \code{list\_recent\_transactions}.
  \item Handmatig: 47 minuten, 4 van de 5 gevonden (V5 tool poisoning gemist), 0 false positives, dus precisie 100\% (4/4).
\end{itemize}

\section{Slide 10. Na indiening: een echte server}
\textbf{Op de deck:} label ``werk na indiening''. Zowe MCP-server (mainframe, mock-modus). Streamable HTTP, 62 tools. Ongeveer 249 bevindingen. ``Precisie stort in: van 92,3\% op Goat naar bijna nul hier.''

\textbf{Achtergrond (niet in de thesis, sluit aan op Validatie, Beperkingen):}
\begin{itemize}
  \item Dit is werk na de indiening, bedoeld om de externe-validiteitsbeperking te toetsen die je zelf benoemt. De thesis stelt al dat de Goat-recall een bovengrens is en de benchmark deels circulair.
  \item Zowe is de offici\"ele open-source mainframe-MCP-stack, gedraaid in mock-modus zonder z/OS-backend en gescand met je bestaande adapter (uitgebreid met SSE-parsing).
\end{itemize}

\section{Slide 11. De findings volgen de policy}
\textbf{Op de deck (tier-tabel):} read(-strict) 31/106/36. update 49/196/62. delete 55/227/74. full 62/249/83. ``De bevindingen schalen met de operator-policy, niet met de beveiligingsstaat.''

\textbf{Achtergrond:}
\begin{itemize}
  \item De capability-tier is het echte autorisatiemechanisme van Zowe. Hoe opener de tier, hoe meer tools, hoe meer findings, lineair. Zelfs maximaal dichtgezet blijven er 36 HIGH-meldingen over op read-only tools.
  \item Oorzaak: jouw heuristieken (SAST-003 en FUZZ-003) zoeken autorisatie in het tool-schema, terwijl Zowe autoriseert op sessie- en transportniveau.
\end{itemize}

\section{Slide 12. Kritische reflectie}
\textbf{Op de deck:}
\begin{itemize}
  \item Circulaire benchmark: Goat is ontworpen voor de regels, dus 100\% recall is best-case.
  \item n=1 reviewer: de handmatige vergelijking is illustratief, niet statistisch valide.
  \item Externe validiteit: hoge recall, lage precisie buiten de benchmark.
\end{itemize}

\textbf{Achtergrond (Validatie, Beperkingen):}
\begin{itemize}
  \item Detectieregels en ground truth zijn niet onafhankelijk, dus 100\% recall generaliseert niet.
  \item De handmatige procedure is door \'e\'en reviewer in een gecontroleerde omgeving uitgevoerd, wat de statistische zeggingskracht beperkt. De thesis noemt dit bewust illustratief.
  \item Het reproduceerbaarheidsverschil is inherent (\code{chess2004static}). MCPTox (\code{wang2025mcptox}) wordt genoemd als methodologische referentie.
\end{itemize}

\section{Slide 13. Aanbevelingen}
\textbf{Op de deck:}
\begin{itemize}
  \item Screeningslaag in CI/CD, geen vervanging van een audit.
  \item Dwing autorisatie af op elke gevoelige tool.
  \item Behandel tool-outputs als onvertrouwde invoer.
  \item Elke run vraagt een menselijke triage.
\end{itemize}

\textbf{Achtergrond (Conclusie, Aanbevelingen):}
\begin{itemize}
  \item ``Het framework vervangt geen security audit. Het verlaagt de drempel om die gericht te voeren.''
  \item V2 en V3 (ontbrekende autorisatie) zijn het hoogste Fintech-risico, dus authenticatie afdwingen, ongeacht of de aanroeper een mens of een agent is.
  \item Tool poisoning in \code{fetch\_market\_news} betekent: outputs sanitiseren zoals inputs.
  \item Een scan onder de 30 seconden maakt CI/CD-integratie haalbaar. Docker verlaagt de drempel.
\end{itemize}

\section{Slide 14. Conclusie}
\textbf{Op de deck:} haalbaar en snel, met hoge recall op de doelklasse. De grens, de precisie buiten de benchmark, is even eerlijk in kaart gebracht. Toekomstig werk: MCPTox (45 re\"ele servers).

\textbf{Achtergrond (Conclusie, Kritische reflectie en Toekomstig onderzoek):}
\begin{itemize}
  \item De hoofdvraag is beantwoord: geautomatiseerde MCP-assessment is realiseerbaar en sneller dan handmatig, bij vergelijkbare of betere recall.
  \item De externe validiteit is expliciet beperkt. Goat representeert niet de gemiddelde productieomgeving.
  \item Drie richtingen voor verder werk: meer en diverse servers (MCPTox, 45 servers en 353 tools, \code{wang2025mcptox}), CI/CD-integratie, en semantische detectie van prompt injection.
\end{itemize}

\section{Back-up-slides}
\begin{itemize}
  \item OWASP MCP Top 10 mapping (correctie op de traceability-tabel die nog LLM07, LLM01 en LLM09 gebruikt): MCP03 tool poisoning, MCP07 autorisatie, MCP02 privilege escalation, Intent Flow Subversion voor prompt injection.
  \item Zowe-adapter en transport.
  \item Scope en beperkingen (geen geauthenticeerde scanning, enkel HTTP en SSE, point-in-time).
\end{itemize}

\end{document}
